% Source: https://artofproblemsolving.com/community/c3h1645439_weird_infinite_series
% Stable identity: aops-topic-1645439-q01
The partial sums are
\[
s_1=a,\quad s_2=a-a=0,\quad s_3=a,\quad s_4=0,\ \dots,
\]
that is, $s_n=a$ for odd $n$ and $s_n=0$ for even $n$.

If $a=0$, every partial sum is $0$, so the series converges with sum $0$.

If $a\ne0$, the partial sums keep jumping between the two different values $a$ and $0$. A limit $L$ would have to be within $\dfrac{|a|}{3}$ of both $a$ and $0$ for all large $n$, which is impossible because $a$ and $0$ are $|a|$ apart. So the series diverges.

Answer: the series converges only for $\boxed{a=0}$, and then its sum is $0$. (Grouping the terms as $(a-a)+(a-a)+\cdots=0$ or as $a-(a-a)-(a-a)-\cdots=a$ gives different ``values'', which shows that brackets cannot be inserted freely in a divergent series. The value $\dfrac {a}{2}$ sometimes quoted is the Cesaro average of the partial sums, not an ordinary sum.)
